% ---------------------------------------------------------
% Breadcrumb
% ---------------------------------------------------------
\begin{tcolorbox}[
  colback=gray!6,
  colframe=gray!40,
  arc=2pt,
  left=8pt, right=8pt, top=6pt, bottom=6pt,
  title={\small\textbf{Where You Are in the Journey}},
  fonttitle=\small\bfseries
]
\begin{center}
\small
Logic, Sets, Proof Techniques
$\;\to\;$ \textbf{Euclidean Geometry}
$\;\to\;$ Analytical Geometry
$\;\to\;$ Non-Euclidean Geometry
$\;\to\;$ $\cdots$
\end{center}

\medskip
\noindent\textbf{How we got here.}
Propositional logic and formal proof gave us the tools to reason
rigorously.
Euclidean geometry is the original domain in which those tools
were deployed: Euclid's \textit{Elements} is the first
sustained axiomatic system in history, and its influence on the
ideal of mathematical proof persists to the present day.

\medskip
\noindent\textbf{What this chapter builds.}
Euclid's five postulates and the synthetic method; congruence
and similarity of triangles; the Pythagorean theorem and its
converses; circle theorems; constructions with compass and
straightedge; area and the method of exhaustion; and the
logical status of the parallel postulate as a genuinely
independent axiom.

\medskip
\noindent\textbf{Where this leads.}
Replacing the parallel postulate with its negation yields
non-Euclidean geometry.
Translating synthetic statements into coordinate equations
yields analytical geometry.
The independence of the parallel postulate, demonstrated
by models of hyperbolic geometry, is also one of the
earliest instances of the model-theoretic reasoning developed
in Volume~I.
\end{tcolorbox}
\vspace{1em}

% ---------------------------------------------------------
% Structural Role
% ---------------------------------------------------------
\begin{tcolorbox}[
  colback=gray!6,
  colframe=gray!40,
  arc=2pt,
  left=8pt, right=8pt, top=6pt, bottom=6pt,
  title={\small\textbf{Structural Role: Classical Geometry — The Axiomatic Template}},
  fonttitle=\small\bfseries
]
Euclidean geometry is the historical and logical origin of the
axiomatic method.

Euclid's program --- state your axioms, define your objects,
derive everything else --- is the template that every other
chapter in this project follows.
Working through the \textit{Elements} therefore serves a dual
purpose: it builds geometric intuition and it makes the
axiomatic method visceral and familiar.

The parallel postulate occupies a special place.
For two millennia, mathematicians attempted to derive it from
the other four.
Its independence, established by consistent non-Euclidean
models, is one of the great conceptual milestones of
mathematics and a direct precursor to the model-theoretic
ideas in Volume~I.
\end{tcolorbox}
\vspace{1em}

% ---------------------------------------------------------
% Structural Roadmap
% ---------------------------------------------------------
\subsection*{Structural Roadmap}

\begin{center}
\textbf{Definitions $\longrightarrow$ Main Theorems
$\longrightarrow$ Consequences and Structural Insight}
\end{center}

The global progression is:

\begin{enumerate}
  \item Euclid's postulates and the synthetic method
  \item Congruence criteria: SAS, ASA, SSS
  \item The Pythagorean theorem and its converse
  \item Similarity and proportionality theorems
  \item Circle theorems: inscribed angles, tangents, power of a point
  \item Compass-and-straightedge constructions and impossibility results
  \item Area via decomposition and the method of exhaustion
  \item Hilbert's axioms: a modern repair of Euclid's foundations
  \item The logical status of the parallel postulate
\end{enumerate}

\begin{remark*}[Primary sources]
Euclid, \textit{Elements} (Heath translation);
Hartshorne, \textit{Geometry: Euclid and Beyond};
Moise, \textit{Elementary Geometry from an Advanced Standpoint}.
\end{remark*}

% ---------------------------------------------------------
% Status
% ---------------------------------------------------------
\vspace{1em}
\begin{tcolorbox}[
  colback=gray!6, colframe=gray!40, arc=2pt,
  left=8pt, right=8pt, top=6pt, bottom=6pt,
  title={\small\textbf{Status: Planned}},
  fonttitle=\small\bfseries
]
\begin{center}\Large\bfseries Coming Soon\end{center}
\vspace{6pt}
\noindent Notes, proofs, and exercises will appear here in a future revision.
\end{tcolorbox}

% Future sections (uncomment as content is written):
% \input{volume-v/geometry/classical-geometry/notes/notes-euclid-postulates}
% \input{volume-v/geometry/classical-geometry/notes/notes-congruence}
% \input{volume-v/geometry/classical-geometry/notes/notes-pythagorean}
% \input{volume-v/geometry/classical-geometry/notes/notes-circles}
% \input{volume-v/geometry/classical-geometry/notes/notes-constructions}
% \input{volume-v/geometry/classical-geometry/notes/notes-area}
% \input{volume-v/geometry/classical-geometry/notes/notes-hilbert-axioms}
% \input{volume-v/geometry/classical-geometry/notes/notes-parallel-postulate}

\clearpage

\section{Book I Definitions}

The following definitions are Euclid's Book I definitions in the Heath
translation \citep[Book I, Definitions 1--23]{euclid-elements-heath}.

\begin{definition}[Point]
\label{def:euclid-point}
A point is that which has no part.
\end{definition}
\begin{lra-not-visible}
\NoLocalDependencies
\end{lra-not-visible}

\begin{definition}[Line]
\label{def:euclid-line}
A line is breadthless length.
\end{definition}
\begin{lra-not-visible}
\NoLocalDependencies
\end{lra-not-visible}

\begin{definition}[Ends of a line]
\label{def:euclid-ends-of-a-line}
The ends of a line are points.
\end{definition}
\begin{lra-not-visible}
\begin{dependencies}
\begin{itemize}
  \item \hyperref[def:euclid-line]{Line}
  \item \hyperref[def:euclid-point]{Point}
\end{itemize}
\end{dependencies}
\end{lra-not-visible}

\begin{definition}[Straight line]
\label{def:euclid-straight-line}
A straight line is a line which lies evenly with the points on itself.
\end{definition}
\begin{lra-not-visible}
\begin{dependencies}
\begin{itemize}
  \item \hyperref[def:euclid-line]{Line}
  \item \hyperref[def:euclid-point]{Point}
\end{itemize}
\end{dependencies}
\end{lra-not-visible}

\begin{definition}[Surface]
\label{def:euclid-surface}
A surface is that which has length and breadth only.
\end{definition}
\begin{lra-not-visible}
\NoLocalDependencies
\end{lra-not-visible}

\begin{definition}[Edges of a surface]
\label{def:euclid-edges-of-a-surface}
The edges of a surface are lines.
\end{definition}
\begin{lra-not-visible}
\begin{dependencies}
\begin{itemize}
  \item \hyperref[def:euclid-line]{Line}
  \item \hyperref[def:euclid-surface]{Surface}
\end{itemize}
\end{dependencies}
\end{lra-not-visible}

\begin{definition}[Plane surface]
\label{def:euclid-plane-surface}
A plane surface is a surface which lies evenly with the straight lines on itself.
\end{definition}
\begin{lra-not-visible}
\begin{dependencies}
\begin{itemize}
  \item \hyperref[def:euclid-straight-line]{Straight line}
  \item \hyperref[def:euclid-surface]{Surface}
\end{itemize}
\end{dependencies}
\end{lra-not-visible}

\begin{definition}[Plane angle]
\label{def:euclid-plane-angle}
A plane angle is the inclination to one another of two lines in a plane which
meet one another and do not lie in a straight line.
\end{definition}
\begin{lra-not-visible}
\begin{dependencies}
\begin{itemize}
  \item \hyperref[def:euclid-line]{Line}
  \item \hyperref[def:euclid-plane-surface]{Plane surface}
  \item \hyperref[def:euclid-straight-line]{Straight line}
\end{itemize}
\end{dependencies}
\end{lra-not-visible}

\begin{definition}[Rectilinear angle]
\label{def:euclid-rectilinear-angle}
And when the lines containing the angle are straight, the angle is called
rectilinear.
\end{definition}
\begin{lra-not-visible}
\begin{dependencies}
\begin{itemize}
  \item \hyperref[def:euclid-plane-angle]{Plane angle}
  \item \hyperref[def:euclid-straight-line]{Straight line}
\end{itemize}
\end{dependencies}
\end{lra-not-visible}

\begin{definition}[Right angle and perpendicular]
\label{def:euclid-right-angle-and-perpendicular}
When a straight line standing on a straight line makes the adjacent angles equal
to one another, each of the equal angles is right, and the straight line
standing on the other is called a perpendicular to that on which it stands.
\end{definition}
\begin{lra-not-visible}
\begin{dependencies}
\begin{itemize}
  \item \hyperref[def:euclid-plane-angle]{Plane angle}
  \item \hyperref[def:euclid-straight-line]{Straight line}
\end{itemize}
\end{dependencies}
\end{lra-not-visible}

\begin{definition}[Obtuse angle]
\label{def:euclid-obtuse-angle}
An obtuse angle is an angle greater than a right angle.
\end{definition}
\begin{lra-not-visible}
\begin{dependencies}
\begin{itemize}
  \item \hyperref[def:euclid-plane-angle]{Plane angle}
  \item \hyperref[def:euclid-right-angle-and-perpendicular]{Right angle and perpendicular}
\end{itemize}
\end{dependencies}
\end{lra-not-visible}

\begin{definition}[Acute angle]
\label{def:euclid-acute-angle}
An acute angle is an angle less than a right angle.
\end{definition}
\begin{lra-not-visible}
\begin{dependencies}
\begin{itemize}
  \item \hyperref[def:euclid-plane-angle]{Plane angle}
  \item \hyperref[def:euclid-right-angle-and-perpendicular]{Right angle and perpendicular}
\end{itemize}
\end{dependencies}
\end{lra-not-visible}

\begin{definition}[Boundary]
\label{def:euclid-boundary}
A boundary is that which is an extremity of anything.
\end{definition}
\begin{lra-not-visible}
\NoLocalDependencies
\end{lra-not-visible}

\begin{definition}[Figure]
\label{def:euclid-figure}
A figure is that which is contained by any boundary or boundaries.
\end{definition}
\begin{lra-not-visible}
\begin{dependencies}
\begin{itemize}
  \item \hyperref[def:euclid-boundary]{Boundary}
\end{itemize}
\end{dependencies}
\end{lra-not-visible}

\begin{definition}[Circle]
\label{def:euclid-circle}
A circle is a plane figure contained by one line such that all the straight
lines falling upon it from one point among those lying within the figure equal
one another.
\end{definition}
\begin{lra-not-visible}
\begin{dependencies}
\begin{itemize}
  \item \hyperref[def:euclid-figure]{Figure}
  \item \hyperref[def:euclid-line]{Line}
  \item \hyperref[def:euclid-plane-surface]{Plane surface}
  \item \hyperref[def:euclid-point]{Point}
  \item \hyperref[def:euclid-straight-line]{Straight line}
\end{itemize}
\end{dependencies}
\end{lra-not-visible}

\begin{definition}[Center of a circle]
\label{def:euclid-center-of-a-circle}
And the point is called the center of the circle.
\end{definition}
\begin{lra-not-visible}
\begin{dependencies}
\begin{itemize}
  \item \hyperref[def:euclid-circle]{Circle}
  \item \hyperref[def:euclid-point]{Point}
\end{itemize}
\end{dependencies}
\end{lra-not-visible}

\begin{definition}[Diameter]
\label{def:euclid-diameter}
A diameter of the circle is any straight line drawn through the center and
terminated in both directions by the circumference of the circle, and such a
straight line also bisects the circle.
\end{definition}
\begin{lra-not-visible}
\begin{dependencies}
\begin{itemize}
  \item \hyperref[def:euclid-center-of-a-circle]{Center of a circle}
  \item \hyperref[def:euclid-circle]{Circle}
  \item \hyperref[def:euclid-straight-line]{Straight line}
\end{itemize}
\end{dependencies}
\end{lra-not-visible}

\begin{definition}[Semicircle]
\label{def:euclid-semicircle}
A semicircle is the figure contained by the diameter and the circumference cut
off by it. And the center of the semicircle is the same as that of the circle.
\end{definition}
\begin{lra-not-visible}
\begin{dependencies}
\begin{itemize}
  \item \hyperref[def:euclid-center-of-a-circle]{Center of a circle}
  \item \hyperref[def:euclid-circle]{Circle}
  \item \hyperref[def:euclid-diameter]{Diameter}
  \item \hyperref[def:euclid-figure]{Figure}
\end{itemize}
\end{dependencies}
\end{lra-not-visible}

\begin{definition}[Rectilinear figures]
\label{def:euclid-rectilinear-figures}
Rectilinear figures are those which are contained by straight lines, trilateral
figures being those contained by three, quadrilateral those contained by four,
and multilateral those contained by more than four straight lines.
\end{definition}
\begin{lra-not-visible}
\begin{dependencies}
\begin{itemize}
  \item \hyperref[def:euclid-figure]{Figure}
  \item \hyperref[def:euclid-straight-line]{Straight line}
\end{itemize}
\end{dependencies}
\end{lra-not-visible}

\begin{definition}[Equilateral, isosceles, and scalene triangles]
\label{def:euclid-triangle-side-classes}
Of trilateral figures, an equilateral triangle is that which has its three
sides equal, an isosceles triangle that which has two of its sides alone equal,
and a scalene triangle that which has its three sides unequal.
\end{definition}
\begin{lra-not-visible}
\begin{dependencies}
\begin{itemize}
  \item \hyperref[def:euclid-rectilinear-figures]{Rectilinear figures}
\end{itemize}
\end{dependencies}
\end{lra-not-visible}

\begin{definition}[Right-angled, obtuse-angled, and acute-angled triangles]
\label{def:euclid-triangle-angle-classes}
Further, of trilateral figures, a right-angled triangle is that which has a
right angle, an obtuse-angled triangle that which has an obtuse angle, and an
acute-angled triangle that which has its three angles acute.
\end{definition}
\begin{lra-not-visible}
\begin{dependencies}
\begin{itemize}
  \item \hyperref[def:euclid-acute-angle]{Acute angle}
  \item \hyperref[def:euclid-obtuse-angle]{Obtuse angle}
  \item \hyperref[def:euclid-rectilinear-figures]{Rectilinear figures}
  \item \hyperref[def:euclid-right-angle-and-perpendicular]{Right angle and perpendicular}
\end{itemize}
\end{dependencies}
\end{lra-not-visible}

\begin{definition}[Quadrilateral classes]
\label{def:euclid-quadrilateral-classes}
Of quadrilateral figures, a square is that which is both equilateral and
right-angled; an oblong that which is right-angled but not equilateral; a
rhombus that which is equilateral but not right-angled; and a rhomboid that
which has its opposite sides and angles equal to one another but is neither
equilateral nor right-angled. And let quadrilaterals other than these be called
trapezia.
\end{definition}
\begin{lra-not-visible}
\begin{dependencies}
\begin{itemize}
  \item \hyperref[def:euclid-rectilinear-angle]{Rectilinear angle}
  \item \hyperref[def:euclid-rectilinear-figures]{Rectilinear figures}
  \item \hyperref[def:euclid-right-angle-and-perpendicular]{Right angle and perpendicular}
\end{itemize}
\end{dependencies}
\end{lra-not-visible}

\begin{definition}[Parallel straight lines]
\label{def:euclid-parallel-straight-lines}
Parallel straight lines are straight lines which, being in the same plane and
being produced indefinitely in both directions, do not meet one another in
either direction.
\end{definition}
\begin{lra-not-visible}
\begin{dependencies}
\begin{itemize}
  \item \hyperref[def:euclid-plane-surface]{Plane surface}
  \item \hyperref[def:euclid-straight-line]{Straight line}
\end{itemize}
\end{dependencies}
\end{lra-not-visible}
